\subsection{现代架构与 Agent 术语补充索引}
\label{intro:24b-modern-lexicon}

这 33 个词条仅补充基础 100 词未覆盖的架构与系统接口，
按\textbf{解决的问题}分组，链接回正文。
例如 RLHF 与 MCP 已在基础词条定义，\textbf{这里不重复列词条}；
英文 A--Z 别名页仅用于查找，不另维护定义。

\subsubsection*{模型输入、目标与训练}
\begin{itemize}
\item \textbf{任务契约（Task Contract）}：事先约定输入、目标、张量尺寸和结果的含义。 \termref{concept:task-contract}
\item \textbf{FFN 与输出头}：前者改变 Token 的隐藏特征，后者映射到词表 logits。 \termref{concept:ffn}
\item \textbf{Teacher Forcing}：训练时以已知真实序列作为条件，一次预测多个后续位置。 \termref{concept:teacher-forcing}
\item \textbf{BPE}：在语料中反复合并高频相邻 Token 对，建立子词单元。 \termref{concept:bpe}
\item \textbf{标签错位（Shifted Targets）}：每个输入位置的目标为下一位置 Token。 \termref{concept:shifted-targets}
\item \textbf{损失掩码}：选择哪些位置参与损失统计，区别于因果注意力掩码。 \termref{concept:loss-mask}
\item \textbf{SFT}：利用示范回答继续训练模型。 \termref{concept:sft}
\item \textbf{RLVR}：利用可程序化核验的任务奖励进行学习。 \termref{concept:rlvr}
\item \textbf{GRPO}：以组内相对奖励构造策略更新信号；还涉及归一化和裁剪等训练步骤。 \termref{concept:grpo}
\item \textbf{Warmup（学习率升温）}：训练前几步由小到大调度学习率，控制初期更新步长。 \termref{concept:lr-warmup}
\item \textbf{AdamW}：保留 Adam 的自适应更新并独立施加权重衰减。 \termref{concept:adamw}
\item \textbf{CoT}：可能要求模型输出中间步骤的提示形式，不等同于直接观察内部计算。 \termref{concept:cot-prompt}
\end{itemize}

\subsubsection*{FFN、位置与归一化}
\begin{itemize}
\item \textbf{RoPE}：对 Q/K 中的分量对做位置相关旋转。 \termref{concept:rope}
\item \textbf{RMSNorm}：按特征的均方根控制数值尺度。 \termref{concept:rmsnorm}
\item \textbf{Pre-Norm}：在某个计算子层之前进行归一化。 \termref{concept:pre-norm}
\item \textbf{SwiGLU}：使用两个投影分支与逐元素门控的 FFN 变体。 \termref{concept:swiglu}
\item \textbf{MoE}：让不同 Token 按路由选择少数专家网络进行计算。 \termref{concept:moe}
\item \textbf{Router（路由器）}：决定当前 Token 由哪些专家处理。 \termref{concept:router}
\end{itemize}

\subsubsection*{Attention 与长上下文}
\begin{itemize}
\item \textbf{MHA}：每个注意力头使用其对应的 Q/K/V 投影。 \termref{concept:multihead}
\item \textbf{MQA}：让多个 Query 头共享一组 K/V 头。 \termref{concept:mqa}
\item \textbf{GQA}：让一组 Query 头共享 K/V 头，降低缓存规模。 \termref{concept:gqa}
\item \textbf{MLA}：通过潜在表示等机制尝试压缩注意力缓存。 \termref{concept:mla}
\item \textbf{KV Cache}：保存过去 Token 的 Key/Value 表示供自回归解码复用。 \termref{concept:kv-cache}
\item \textbf{滑动窗口 Attention}：每个位置只读取一个有限的可见前缀窗口。 \termref{concept:sliding-window-attention}
\item \textbf{稀疏 Attention}：只计算选择的 Token 对，不要求全部两两比较。 \termref{concept:sparse-attention}
\item \textbf{线性 Attention}：通过不同运算结构改善长序列计算或状态扩展性，通常不等同于原 Softmax 公式。 \termref{concept:linear-attention}
\item \textbf{外积（Outer product）}：将列形式 Key 与行形式 Value 相乘，写成关联矩阵。 \termref{concept:outer-product}
\item \textbf{串扰（Interference）}：多个不正交 Key 的读取贡献叠加，污染目标 Value。 \termref{concept:memory-interference}
\item \textbf{Delta 规则}：先读旧状态，用目标与旧输出的残差修正写入。 \termref{concept:delta-rule}
\end{itemize}

\subsubsection*{Agent 系统中的边界}
\begin{itemize}
\item \textbf{Tool Calling}：模型提出结构化请求，系统在检查后调用外部工具。 \termref{concept:tool-calling}
\item \textbf{Workflow}：由固定步骤或条件分支编排的可重复任务流程。 \termref{concept:workflow}
\item \textbf{Skill}：把一类任务的说明、约束和脚本包装为可复用能力。 \termref{concept:skill}
\item \textbf{Subagent 与 A2A}：前者为子任务分配独立执行单元；后者规范独立 Agent 间通信。 \termref{concept:subagent}
\end{itemize}

\begin{conclusion}
遇到新缩写，先问它改变了\textbf{数据、模型层、训练信号、
计算复杂度、缓存，还是系统权限}。
明确所属层级，再从这里点击回正文，远比背名称更有用。
\end{conclusion}
